\clearpage
\begin{examplebox}
\paragraph{例1：估算 $e^{0.1}$，并保证四位小数}
识别 $x=0.1\in[0,1]$，代入指数分式：
\[
R_E(0.1)=\frac{1261}{1141}=1.10517090\ldots.
\]
由已证误差界，
\[
0<e^{0.1}-R_E(0.1)<\frac{3(0.1)^5}{665}<4.52\times10^{-8}.
\]
因而 $1.10517090<e^{0.1}<1.10517096$，整个区间都落在舍入到 $1.1052$ 的范围 $(1.10515,1.10525)$ 内，故\textbf{四位小数为 $1.1052$}。
回验依靠误差区间，不必预先知道计算器中的真实值。
变式：由 $R_E(-x)=1/R_E(x)$，估算 $e^{-x}$ 时可把指数分式的分子、分母互换；$x>0$ 时，这个分式高估 $e^{-x}$。

\paragraph{例2：估算 $\ln1.1$，并判断高估还是低估}
这里 $1.1=1+0.1$，因此 $x=0.1$，而不是 $1.1$。
\[
R_L(0.1)=\frac{63}{661}=0.09531013\ldots,
\qquad
0<\ln1.1-R_L(0.1)<\frac{10^{-5}}{180}.
\]
所以 $0.09531013<\ln1.1<0.09531020$，分式是低估，\textbf{四位小数为 $0.0953$}。
回验也可看面积：矩形总面积接近曲线下面积；其严格偏差方向由方法一保证，不能仅凭肉眼判断。

\newpage
\paragraph{例3：$\tan0.2$ 的近似与独立夹逼}
$0.2$ 是弧度数，且位于零点附近。代入得
\[
R_T(0.2)=\frac{0.6}{2.96}=\frac{15}{74}=0.2027027\ldots.
\]
这提示四位小数可能是 $0.2027$，但代入结果本身还不是精度证明。
独立回验用 $x=0.2$ 时的正弦、余弦夹逼：
\[
s_-=x-\frac{x^3}{6}<\sin x<s_+=x-\frac{x^3}{6}+\frac{x^5}{120},
\]
\[
0<c_-=1-\frac{x^2}{2}<\cos x<c_+=1-\frac{x^2}{2}+\frac{x^4}{24}.
\]
这些方向由拉格朗日余项及 $0<\xi<x<\pi/2$ 时 $\sin\xi>0$ 确定。
分子、分母均为正，故
\[
0.20270<\frac{s_-}{c_+}<\tan0.2<\frac{s_+}{c_-}<0.20273.
\]
因此四位小数确实为 $0.2027$。注意：这不等于分式给出了更多位的准确数字。
\begin{center}
\begin{tikzpicture}[font=\small]
\draw[black!30,rounded corners=2pt] (0,0) rectangle (14.4,8.2);
\node[anchor=west,font=\small\bfseries] at (0.5,7.72) {三角函数：绝对误差随自变量变化};
\begin{axis}[width=11.2cm,height=6.1cm,at={(2.05cm,1.25cm)},anchor=south west,
axis lines=left,xlabel={$x$},tick label style={font=\small},label style={font=\small},
legend style={font=\small,draw=none,fill=white,at={(0.03,0.97)},anchor=north west},
legend cell align=left,unbounded coords=jump,ymode=log,xmin=0.1,xmax=1,ymin=1e-10,ymax=0.1,ylabel={绝对误差},restrict x to domain=0.1:1,legend columns=3,legend style={font=\small,draw=none,fill=none,at={(0.5,1.15)},anchor=south,column sep=3pt}]
\addplot[blue,very thick] table[x=x,y expr={abs(\thisrow{es})}] {assets/tikz/F003_samples.dat};
\addlegendentry{$|\sin x-R_S(x)|$}
\addplot[orange!85!black,thick,dashed] table[x=x,y expr={abs(\thisrow{ec})}] {assets/tikz/F003_samples.dat};
\addlegendentry{$|\cos x-R_C(x)|$}
\addplot[green!40!black,thick,densely dotted] table[x=x,y expr={abs(\thisrow{et})}] {assets/tikz/F003_samples.dat};
\addlegendentry{$|\tan x-R_T(x)|$}

\end{axis}
\end{tikzpicture}
\end{center}

图示 $0.1\le x\le1$ 的误差，三条曲线来自原函数与各自分式的实际差值；在 $x=0.2$ 处，正切分式误差约 $7.33\times10^{-6}$，不能把它输出的所有小数位都当作有效数字。
\end{examplebox}
